\section{Exact finite certificate and claim ledger}
\label{sec:tpc43-certificate}

The companion script
\begin{center}
\texttt{experiments/tpc43\_certificate.py}
\end{center}
uses only Python's standard library and exact integer, rational, and
Gaussian--integer arithmetic.  It performs no random sampling.  Its
finite regressions verify:
\begin{enumerate}[label=(\roman*)]
 \item the squarefree--kernel character identity
 \eqref{eq:tpc43-fused-Walsh-character};
 \item the divisor reconstruction and multiplicity bound in
 \cref{lem:tpc43-kernel-fiber-bound};
 \item zero--alias fused--label injectivity on exhaustive admissible toy
 packets;
 \item exact Walsh second moments and the diagonal identity
 \eqref{eq:tpc43-exact-zero-expectation};
 \item the four--to--one identity
 \eqref{eq:tpc43-literal-four-to-one};
 \item the rough--kernel filtration
 \eqref{eq:tpc43-rough-filtration};
 \item exact Walsh Parseval for the grouped spectrum; and
 \item the sharp one--corner extremizer in
 \cref{thm:tpc43-one-corner-extremizer}.
\end{enumerate}
The script also checks the endpoint exponent ledger and scans its own
syntax tree to exclude floating--point literals, random imports, true
division, and optimization--sensitive \texttt{assert} statements.
The frozen run performs (43{,}877) exact checks.  Its canonical
certificate digest is
\begin{center}
\texttt{d7d2b07107a6da1bd242ef9dbd8c02de7daea24dfdfbf9b5f31178cacf834d01}.
\end{center}
Normal and optimized Python executions produce the same JSON file, whose
SHA--256 digest is
\begin{center}
\texttt{f46e5b5ef49cc5f12f92ba0f702938a3fa2367b30c694e01f93fd14d30bd628a}.
\end{center}

Finite computation certifies the algebraic identities and finite
combinatorics only.  The asymptotic divisor bound, Bonami--Beckner
inequality, inherited TPC diagonal estimates, and Borel--Cantelli lemma
are mathematical inputs proved or cited in the text.  The certificate
does not test a physical M\"obius asymptotic, affine Chowla cancellation,
prime pairs, or twin primes.
